\section{Summary}
\label{sec:summary}

This project developed and evaluated a stateful Slurm Agent that enables natural-language interaction for HPC cluster monitoring and control. The final system combines a FastAPI backend, OpenAI Agents SDK orchestration, an MCP-based tool layer, and a React + Vite frontend with integrated chat and evaluation modes.

\subsection{Research Contributions}

This work contributes the following practical outcomes:

\textbf{1. Scenario-Grid Evaluation for Agentic Slurm Workflows.} The project introduces a dataset-driven evaluation design (455 curated cases plus a separate 40-case paraphrase split) that measures architecture-level behavior: tool usage, routing, HITL safety, state transitions, and response quality.

\textbf{2. Stateful MCP-Centered Integration.} Slurm interactions are exposed through MCP tools and exercised in deterministic mock scenarios with per-test reset, enabling reproducible behavioral testing.

\textbf{3. Multi-Agent Reasoning and Safety Flow.} Observer/operator delegation and confirmation workflows are integrated into a single runtime, including explicit pending-action management.

\textbf{4. Structured Streaming Interface.} The backend emits structured SSE deltas (reasoning, tool output, content, pending actions, chart artifacts), allowing richer frontend rendering and traceability.

\textbf{5. Integrated Evaluation UI.} Evaluation is available as a first-class frontend mode, reducing friction between implementation changes and benchmark validation.

\subsection{Achievement Against Project Goals}

\textbf{Goal 1 (Natural language access to Slurm):} Achieved for read-oriented and account-oriented workflows, with strong category-level performance.

\textbf{Goal 2 (Multi-agent orchestration):} Achieved. Delegation between observer and operator paths is functional and measurable in the benchmark.

\textbf{Goal 3 (MCP standardization):} Achieved. Tool interfaces are separated from agent reasoning and can be extended independently.

\textbf{Goal 4 (Safety for dangerous operations):} Partially achieved. HITL flow is implemented and active, but benchmark safety metrics indicate additional hardening is needed.

\textbf{Goal 5 (Realtime UX and web interface):} Achieved with structured streaming in the React frontend and integrated evaluation controls.

\subsection{Key Quantitative Outcomes}

\begin{table}[H]
    \centering
    \caption{Final Snapshot Highlights}
    \label{tab:summary-highlights}
    \begin{tabular}{|l|c|}
        \hline
        \textbf{Metric} & \textbf{Value} \\
        \hline
        Dataset size & 455 curated tests \\
        Pass rate ($\geq 0.80$) & 93.6\% (426/455) \\
        Average overall score & 0.936 \\
        BAR & 0.908 \\
        SVR & 0.161 \\
        CSR & 0.934 \\
        Average latency & 22.59s \\
        \hline
    \end{tabular}
\end{table}

These results show strong capability on the curated benchmark and acceptable cross-scenario stability. They do not by themselves prove unrestricted robustness to arbitrary Slurm requests; that question is addressed separately through the held-out paraphrase split and future live-cluster validation.

\subsection{Main Findings}

\textbf{Architecture-aware metrics are necessary.} Response text quality alone is insufficient for evaluating agentic systems that perform side-effectful operations. Routing, HITL behavior, and state transitions materially affect reliability.

\textbf{Deterministic state reset improves evaluation quality.} Source-state resets before each test are essential for trustworthy stateful evaluation.

\textbf{Frontend-integrated evaluation improves iteration speed.} Embedding the evaluation workflow in the main UI shortens the feedback loop for feature changes.

\textbf{Safety remains a hard problem.} Even with confirmation mechanisms, mismatch rates in safety-related categories show that stronger guarantees are still required before production use.

\subsection{Source Code Availability}

The complete source code is publicly available:

\begin{center}
    \url{https://github.com/DanhNguyennene/specialized_project_slurm_agent}
\end{center}

The repository contains the backend, MCP server, React frontend, evaluation framework, dataset artifacts, and reproducible result snapshots.

\subsection{TODO for Final Thesis Freeze}

\begin{itemize}
    \item Freeze one final 455-case curated benchmark snapshot under the pinned local-Ollama stack (qwen3.5:9b main, qwen2.5:7b specialist, gpt-oss:20b judge) and update all chapter references to that single run.
    \item Run the 40-case held-out paraphrase split and report it separately from the curated benchmark.
    \item Add one concise reliability paragraph interpreting pass\textasciicircum{}k once repeat-run data is finalized.
\end{itemize}